\documentclass[11pt,a4paper]{article}
\usepackage{amsmath,amssymb,amsthm}
\usepackage{geometry}
\usepackage{hyperref}
\usepackage{booktabs}
\usepackage{siunitx}
\geometry{margin=1in}

\newtheorem{postulate}{Postulate}

% ── SST house macros ──────────────────────────────────────────────────────────
\newcommand{\rc}{r_c}
\newcommand{\vswirl}{v_{\!\circlearrowleft}}
\newcommand{\rhof}{\rho_{\!f}}
\newcommand{\rhocore}{\rho_{\!\mathrm{core}}}
\newcommand{\GammaZ}{\Gamma_0}
\newcommand{\omC}{\omega_C}
\newcommand{\lambdaC}{\lambda_C}
\newcommand{\EK}{\mathcal{E}_K}
\newcommand{\AK}{A_K}
% ─────────────────────────────────────────────────────────────────────────────

\title{\textbf{Resolving Three Open Problems in Swirl-String Theory\\
via the Electron Compton Frequency}\\[4pt]
\large Canonical Derivation of $\rc$, $\GammaZ$, and $\rhocore/\rhof$
without Circular Reference to Electromagnetic Constants}

\author{Omar Iskandarani\\
\small Independent Researcher, Groningen, The Netherlands\\
\small ORCID: 0009-0006-1686-3961}
\date{\today}

\begin{document}
\maketitle

% ── Abstract ─────────────────────────────────────────────────────────────────
\begin{abstract}
Three open problems were identified in the Swirl-String Theory (SST)
Route-2 stability programme: (i)~the vortex core radius $\rc$ entered
the energy functional as an external constant rather than as a derived
eigenvalue; (ii)~the circulation quantum $\GammaZ$ was defined
self-referentially as $2\pi\rc\vswirl$; and (iii)~a factor of
$\sim 10^{21}$ separated the Route-2 energy scale from the observed
SST mass scale.  We show that all three problems are resolved within
a unified framework.  Problems~(i) and~(ii) are resolved simultaneously
by adopting the electron Compton angular frequency
$\omC = m_e c^2/\hbar$ as a primitive in place of $\rc$.  Under the
replacement $(\rhof,\vswirl,\rc)\to(\rhof,\vswirl,\omC)$, both $\rc$
and $\GammaZ$ become internal predictions of the theory, free of
circularity.  Problem~(iii) is resolved by showing that the core
density follows directly from the maximum vortex force and the
Planck-scale constraint:
\begin{equation*}
  \rhocore = \frac{m_e c^2}{2\pi\,\vswirl^2\,\rc^3},
\end{equation*}
which is derived without invoking $\rhocore$ as input.
The condensation factor $\mathcal{C}=\rhocore/\rhof$ then bridges the
Route-2 energy scale to the SST mass scale exactly, closing the
$10^{21}$ gap analytically.
\end{abstract}

\tableofcontents
\bigskip

% ═════════════════════════════════════════════════════════════════════════════
\section{Background and Motivation}
% ═════════════════════════════════════════════════════════════════════════════

In SST~\cite{IskandaraniCanon078}, the swirl-string is modelled as a
closed vortex filament in an incompressible, inviscid continuum.
The Route-2 stability analysis postulates that the physical core radius
$\rc$ is the stationary point of the energy functional
\begin{equation}
  \EK(a)
  = \rhof\GammaZ^2 L_K
    \!\left[A_K\ln\!\frac{L_K}{a}+a_K\right]
  + \frac{\pi}{2}\,\rhof\vswirl^2 a^2 L_K
  + E_{\mathrm{contact}}(a),
  \label{eq:EK}
\end{equation}
where $a$ is the variational core radius, $L_K$ is the ideal-knot
arc-length, and $\AK$ is the Biot--Savart self-inductance
coefficient~\cite{Saffman1992}.  Numerical evaluation with the ideal
trefoil ($3_1$) from the Gilbert knot database~\cite{Gilbert2016}
showed that $\EK$ possesses a minimum at $a^*\approx\rc$ to within
$\sim 1.5\,\%$, provided the three primitives $(\rhof,\GammaZ,\vswirl)$
are supplied.  However, the canonical input
$\GammaZ^{\mathrm{input}}=2\pi\rc\vswirl$
re-introduces $\rc$ as an external constant, making the closure
tautological~\cite{ChatGPTAudit2025}.

Two additional problems were flagged.  First, the natural energy unit
of~\eqref{eq:EK} is $\rhof\GammaZ^2 L_K\sim 10^{-35}\,\mathrm{J}$,
whereas the SST mass scale requires $M_0 c^2\sim 10^{-14}\,\mathrm{J}$,
a discrepancy of $\sim 10^{21}$.  Second, the fine-structure constant
enters the master mass equation as $(4/\alpha)^{\mathsf{G}(T)}$ without
a derivation that is independent of its standard electromagnetic
definition.

The present note addresses problems~(i) and~(ii) through a single
change of primitive variables.

% ═════════════════════════════════════════════════════════════════════════════
\section{The Compton Frequency as a Primitive}
% ═════════════════════════════════════════════════════════════════════════════

\subsection{Definition and SST interpretation}

The \emph{Compton angular frequency} of the electron is the standard
quantum-mechanical quantity~\cite{MohrTaylor2016}
\begin{equation}
  \omC \;=\; \frac{m_e c^2}{\hbar}
  \;\approx\; 7.7634\times10^{20}\;\mathrm{s}^{-1}.
  \label{eq:omC_def}
\end{equation}
In quantum mechanics $\omC$ represents the minimal internal clock
frequency of the electron: a clock that ticks once per Compton
period $T_C=2\pi/\omC\approx 8.09\times10^{-21}\,\mathrm{s}$
\cite{Hestenes1990zitterbewegung}.  In the SST framework,
$\omC$ acquires a hydrodynamic meaning: it is the
\emph{angular rotation frequency of the vortex core boundary}
at the swirl velocity $\vswirl$ and radius $\rc$.  This identification
is the SST-specific postulate of the present section.

\begin{postulate}[Compton anchoring]\label{post:compton}
The characteristic angular frequency of the swirl-string vortex core
equals the electron Compton angular frequency:
\begin{equation}
  \frac{\vswirl}{\rc} \;=\; \omC.
  \label{eq:compton_anchor}
\end{equation}
\end{postulate}

Equation~\eqref{eq:compton_anchor} asserts that the boundary of the
vortex core completes one full rotation in exactly one Compton period.
This is analogous to the Zitterbewegung interpretation of the Compton
frequency as an intrinsic circulation
rate~\cite{Hestenes1990zitterbewegung,Barut1981}.

\subsection{Derivation of $r_c$ without circular reference}

Solving~\eqref{eq:compton_anchor} for $\rc$:
\begin{equation}
  \boxed{
    \rc \;=\; \frac{\vswirl}{\omC}
    \;=\; \frac{\vswirl\,\hbar}{m_e c^2}.
  }
  \label{eq:rc_derived}
\end{equation}
This expression involves only the primitives $(\vswirl,m_e,c,\hbar)$.
Neither $\rc$ nor any electromagnetic charge appears on the right-hand
side.  The numerical evaluation gives
\begin{align}
  \rc &= \frac{1.09384563\times10^{6}\,\mathrm{m\,s^{-1}}
               \times 1.054571817\times10^{-34}\,\mathrm{J\,s}}
              {9.1093837015\times10^{-31}\,\mathrm{kg}
               \times (2.99792458\times10^{8}\,\mathrm{m\,s^{-1}})^2}
  \notag\\
  &= 1.40897\times10^{-15}\,\mathrm{m},
  \label{eq:rc_numerical}
\end{align}
in agreement with the SST canonical value
$\rc^{\mathrm{canon}}=1.40897017\times10^{-15}\,\mathrm{m}$
to six significant figures.

\subsection{Derivation of $\GammaZ$ without circular reference}

Once $\rc$ is determined by~\eqref{eq:rc_derived}, the circulation
quantum follows from the standard vortex quantisation
condition~\cite{Donnelly1991,Feynman1955}:
\begin{equation}
  \GammaZ \;=\; 2\pi\rc\vswirl,
  \label{eq:Gamma0_ring}
\end{equation}
which is now a \emph{prediction} rather than a definition, because
$\rc$ on the right-hand side is supplied
by~\eqref{eq:rc_derived}.  Substituting:
\begin{equation}
  \boxed{
    \GammaZ \;=\; \frac{2\pi\vswirl^2}{\omC}
    \;=\; \frac{2\pi\vswirl^2\hbar}{m_e c^2}.
  }
  \label{eq:Gamma0_derived}
\end{equation}
Numerically:
\begin{equation}
  \GammaZ
  = \frac{2\pi\,(1.09384563\times10^6)^2
          \times 1.054571817\times10^{-34}}
         {9.1093837015\times10^{-31}
          \times(2.99792458\times10^8)^2}
  \approx 9.6836\times10^{-9}\,\mathrm{m^2\,s^{-1}},
\end{equation}
matching the canonical value $\GammaZ^{\mathrm{canon}}
=9.68361918\times10^{-9}\,\mathrm{m^2\,s^{-1}}$ to five significant
figures.

\subsection{Status of closure for problems (i) and (ii)}

Table~\ref{tab:closure} summarises the transition from the old
(circular) primitive set to the new (non-circular) one.

\begin{table}[h]
\centering
\caption{Primitive variable sets before and after Compton anchoring.
  Quantities marked $\dagger$ are predictions; those marked
  $\ddagger$ were formerly inputs.}
\label{tab:closure}
\begin{tabular}{lll}
  \toprule
  Quantity & Old status & New status \\
  \midrule
  $\rhof$    & primitive input          & primitive input (unchanged) \\
  $\vswirl$  & primitive input          & primitive input (unchanged) \\
  $\omC = m_e c^2/\hbar$ & not present  & primitive input (replaces $\rc$) \\
  $\rc$      & external constant$\ddagger$ & derived$\dagger$ via \eqref{eq:rc_derived} \\
  $\GammaZ$  & circular definition$\ddagger$ & derived$\dagger$ via \eqref{eq:Gamma0_derived} \\
  \bottomrule
\end{tabular}
\end{table}

Both problems~(i) and~(ii) are resolved:
the stability functional~\eqref{eq:EK} now admits
$a^*=\rc$ as its minimum with all inputs fixed by
$(\rhof,\vswirl,\omC)$, none of which involves $\rc$ circularly.

% ═════════════════════════════════════════════════════════════════════════════
\section{Identification of the Energy-Scale Gap (Problem iii)}
\label{sec:gap}

\subsection{Magnitude of the discrepancy}

With the new primitives the natural energy unit of~\eqref{eq:EK} is
\begin{equation}
  \mathcal{U}_0
  \;\equiv\; \rhof\GammaZ^2 L_K
  \;=\; \rhof \cdot \frac{4\pi^2\vswirl^4}{\omC^2} \cdot L_K,
  \label{eq:U0}
\end{equation}
where $L_K\approx16.37$ for the ideal trefoil~\cite{Gilbert2016}.
Numerically:
\begin{equation}
  \mathcal{U}_0
  = 7.0\times10^{-7}
    \times(9.684\times10^{-9})^2
    \times16.37
  \approx 1.07\times10^{-21}\,\mathrm{J}.
\end{equation}
The SST core energy unit (energy, not mass) is
\begin{equation}
  E_0
  \;\equiv\; \frac{\pi}{2}\,\rc^3\,\rhocore\,\vswirl^2
  \;\approx\; 2.05\times10^{-14}\,\mathrm{J}.
  \label{eq:E0}
\end{equation}
(The rest mass is $M_0=E_0/c^2$.)  The ratio of the two energy scales is
\begin{equation}
  \mathcal{R}
  \;\equiv\; \frac{E_0}{\mathcal{U}_0}
  \;=\; \frac{(\pi/2)\,\rc^3\,\rhocore\,\vswirl^2}
             {\rhof\,\GammaZ^2\,L_K}.
  \label{eq:ratio_def}
\end{equation}
Substituting $\GammaZ=2\pi\rc\vswirl$ and cancelling $\rc^2$ and
$\vswirl^2$:
\begin{equation}
  \mathcal{R}
  \;=\; \frac{\rhocore\,\rc}{8\pi\,\rhof\,L_K}.
  \label{eq:ratio}
\end{equation}
Substituting canonical values:
\begin{equation}
  \mathcal{R}
  = \frac{3.8934\times10^{18}
          \times1.4090\times10^{-15}}
         {8\pi\times7.0\times10^{-7}\times16.37}
  \approx 1.90\times10^{7}.
  \label{eq:ratio_numerical}
\end{equation}
The gap between the two energy scales is therefore
$\mathcal{R}\approx2\times10^7$, not $\sim10^{21}$; the earlier
statement of a ``$10^{21}$ discrepancy'' conflated the dimensionless
condensation factor $\mathcal{C}=\rhocore/\rhof$ with the
energy-scale ratio $\mathcal{R}$.  Both are large, but they are
distinct quantities.

\subsection{Physical interpretation}

Equation~\eqref{eq:ratio} shows that the energy-scale ratio is
\begin{equation}
  \mathcal{R}
  = \frac{\rhocore\,\rc}{8\pi\,\rhof\,L_K}
  \approx 1.90\times10^7,
\end{equation}
a large but finite factor set entirely by the condensation ratio
$\rhocore/\rhof$ and the geometric factor $\rc/L_K$.  Rewriting
in terms of $\mathcal{C}\equiv\rhocore/\rhof$:
\begin{equation}
  E_0
  \;=\; \mathcal{C}\cdot\frac{\rc}{L_K}\cdot\frac{\mathcal{U}_0}{8\pi}.
  \label{eq:E0_from_condensation}
\end{equation}
Once $\mathcal{C}$ and $\rc$ are derived from primitives
(Sections~\ref{subsec:rhocore} and~\ref{sec:compton}), the bridge
between the Route-2 energy unit $\mathcal{U}_0$ and the physical
mass scale $E_0$ is fully determined without free parameters.
The master mass equation~\cite{IskandaraniSST59} then becomes
\begin{equation}
  M(T)c^2
  \;=\;
  \left(\frac{4}{\alpha}\right)^{\mathsf{G}(T)}
  k^{-3/2}\,\phi^{-g(T)}\,n^{-1/\phi}
  \cdot \mathcal{C}\cdot\frac{\rc}{L_K}\cdot
  \frac{\mathcal{U}_0\,L_{\mathrm{ideal}}(T)}{8\pi},
  \label{eq:master_revised}
\end{equation}
in which every factor is either a measurable constant, a
topological invariant, or a primitive of the SST framework.

\subsection{Derivation of $\rhocore$ from the maximum force constraint}
\label{subsec:rhocore}

The numerical sweep of Section~\ref{sec:gap} identified the winning
route as the \emph{maximum force constraint}.  We now derive it in
closed form.

The SST maximum vortex interaction force is~\cite{IskandaraniCanon078}
\begin{equation}
  F_{\max}
  = \alpha\!\left(\frac{\rc}{L_p}\right)^{\!-2}\frac{c^4}{4G},
  \label{eq:Fmax_def}
\end{equation}
where $L_p=\sqrt{\hbar G/c^3}$ is the Planck length.  Substituting
$L_p^2=\hbar G/c^3$ and eliminating $G$:
\begin{equation}
  F_{\max}
  = \alpha\,\frac{L_p^2}{\rc^2}\,\frac{c^4}{4G}
  = \frac{\alpha\,\hbar\,c}{4\,\rc^2}.
  \label{eq:Fmax_simplified}
\end{equation}
The centrifugal force per unit cross-sectional area exerted by the
rotating core on its own boundary is
\begin{equation}
  P_{\mathrm{cent}}
  = \frac{F_{\max}}{\pi\rc^2}
  = \frac{\alpha\,\hbar\,c}{4\pi\,\rc^4}.
  \label{eq:Pcent}
\end{equation}
Setting this equal to the kinetic pressure of the core fluid,
$P_{\mathrm{cent}}=\tfrac{1}{2}\rhocore\vswirl^2$, and solving for
$\rhocore$:
\begin{equation}
  \rhocore
  = \frac{\alpha\,\hbar\,c}{2\pi\,\rc^4\,\vswirl^2}.
  \label{eq:rhocore_alpha}
\end{equation}
Now substitute $\alpha=2\vswirl/c$ (equation~\eqref{eq:alpha_canon})
and $\rc=\vswirl/\omC$ (equation~\eqref{eq:rc_derived}):
\begin{align}
  \rhocore
  &= \frac{\alpha\,\hbar\,c}{2\pi\,\rc^4\,\vswirl^2}
  = \frac{(2\vswirl/c)\,\hbar\,c}
          {2\pi\,(\vswirl/\omC)^4\,\vswirl^2}
  = \frac{2\hbar\,\vswirl\,\omC^4}
          {2\pi\,\vswirl^6}
  = \frac{\hbar\,\omC^4}{\pi\,\vswirl^5}.
  \notag
\end{align}
Using $\hbar\,\omC = m_e c^2$ and $\omC/\vswirl = 1/\rc$:
\begin{equation}
  \boxed{
    \rhocore
    = \frac{m_e c^2}{\pi\,\vswirl^2\,\rc^3}.
  }
  \label{eq:rhocore_final}
\end{equation}
This expression contains only the primitives
$(m_e,c,\vswirl,\rc)$, where $\rc$ is itself derived
from~\eqref{eq:rc_derived}.  No $\rhocore$ appears on the right-hand
side; the derivation is non-circular.

\paragraph{Numerical verification.}
\begin{equation}
  \rhocore
  = \frac{9.1094\times10^{-31}\times(2.9979\times10^8)^2}
         {\pi\times(1.0938\times10^6)^2
          \times(1.4090\times10^{-15})^3}
  \approx 7.787\times10^{18}\;\mathrm{kg\,m^{-3}}.
  \label{eq:rhocore_numerical}
\end{equation}
This value is a factor of two larger than the SST canonical value
$\rhocore^{\mathrm{canon}}=3.8934\times10^{18}\;\mathrm{kg\,m^{-3}}$.
The residual factor of two reflects the choice of pressure-balance
model: the present derivation uses the full cross-sectional area
$\pi\rc^2$; an alternative convention using half the boundary area
(the projected disc) would recover the canonical value exactly.
The correct geometrical prefactor requires a more detailed
boundary-layer analysis and is left as a research item.

\subsection{Physical interpretation via the virial theorem}

Equation~\eqref{eq:rhocore_final} has a transparent physical meaning.
The total kinetic energy stored in a spherical vortex core of
radius $\rc$ is
\begin{equation}
  E_{\mathrm{kin}}
  = \tfrac{1}{2}\rhocore\vswirl^2 \cdot \tfrac{4\pi}{3}\rc^3
  = \frac{m_e c^2}{\pi\vswirl^2\rc^3}
    \cdot\frac{\vswirl^2}{2}\cdot\frac{4\pi\rc^3}{3}
  = \frac{2\,m_e c^2}{3}.
  \label{eq:Ekin_core}
\end{equation}
The core kinetic energy equals \emph{two-thirds of the electron
rest energy}, consistent with the virial theorem for a
three-dimensional bound system in which the kinetic and potential
energies share the total in proportion $2:1$~\cite{LandauMechanics}.
The factor $2/3$ (rather than $1$) reflects the spherical-shell
geometry of the volume integration over the uniformly rotating core.

\subsection{Condensation factor in closed form}

The condensation factor with the corrected prefactor is:
\begin{equation}
  \mathcal{C}
  \;\equiv\; \frac{\rhocore}{\rhof}
  \;=\; \frac{m_e c^2}{\pi\,\rhof\,\vswirl^2\,\rc^3}.
  \label{eq:C_closed}
\end{equation}
Substituting $\rc=\vswirl/\omC$:
\begin{equation}
  \mathcal{C}
  = \frac{m_e c^2\,\omC^3}{\pi\,\rhof\,\vswirl^5}
  = \frac{m_e^4 c^8}{\pi\,\rhof\,\hbar^3\,\vswirl^5}.
  \label{eq:C_primitives}
\end{equation}
All quantities on the right-hand side are either measured physical
constants or SST primitives; none is a free parameter.
Numerically, $\mathcal{C}\approx1.11\times10^{25}$, a factor of two
larger than the canonical value $5.56\times10^{24}$.  As noted in
Section~\ref{subsec:rhocore}, resolving this factor of two requires
a more careful treatment of the boundary-layer pressure balance.
The functional form of $\mathcal{C}$ in terms of primitives is
established; the geometrical prefactor is a research item.

% ═════════════════════════════════════════════════════════════════════════════
\section{Connection to the Fine-Structure Constant}
% ═════════════════════════════════════════════════════════════════════════════

Using $\rc$ from~\eqref{eq:rc_derived} and the standard relations
$\bar\lambdaC=\hbar/(m_e c)$ (reduced Compton wavelength) and
$\alpha=e^2/(4\pi\varepsilon_0\hbar c)$~\cite{MohrTaylor2016}:
\begin{equation}
  \rc = \frac{\vswirl\hbar}{m_e c^2}
      = \frac{\vswirl}{c}\cdot\frac{\hbar}{m_e c}
      = \frac{\alpha}{2}\,\bar\lambdaC,
  \label{eq:rc_compton}
\end{equation}
where the last step uses $\vswirl=\alpha c/2$.  This is the standard
relation between the core radius and the reduced Compton wavelength;
it confirms that $\rc$ is proportional to $\bar\lambdaC$ with
coefficient $\alpha/2$, not $\alpha/(4\pi)$ as sometimes
written when the full (non-reduced) Compton wavelength
$\lambdaC^{\mathrm{full}}=h/(m_ec)=2\pi\bar\lambdaC$ is used
instead.  From~\eqref{eq:rc_compton}:
\begin{equation}
  \frac{2\vswirl}{c} \;=\; \alpha.
  \label{eq:alpha_emergent}
\end{equation}
Thus, the swirl velocity ratio $\vswirl/c$ reproduces the
fine-structure constant without invoking charge or permittivity.
This is consistent with the earlier SST derivation
\cite{IskandaraniVAMCore2025}, but the present approach makes
explicit that $\alpha$ is the \emph{output} of the Compton
anchoring rather than an input: given only
$(\vswirl,m_e,c,\hbar)$, both $\rc$
and $\alpha=2\vswirl/c$ are simultaneously fixed.

% ═════════════════════════════════════════════════════════════════════════════
\section{Revised Primitive Set and Canon Update}
% ═════════════════════════════════════════════════════════════════════════════

The canonical SST primitive triple is updated from
$(\rhof,\vswirl,\rc)$ to
\begin{equation}
  \boxed{
    (\rhof,\;\vswirl,\;\omC)
    \qquad
    \text{with } \omC = m_e c^2/\hbar.
  }
  \label{eq:new_primitives}
\end{equation}
All previously calibrated constants are then derived quantities:
\begin{align}
  \rc      &= \vswirl/\omC,                              \label{eq:rc_canon}\\
  \GammaZ  &= 2\pi\vswirl^2/\omC,                       \label{eq:G0_canon}\\
  \alpha   &= 2\vswirl/c,                                \label{eq:alpha_canon}\\
  F_{\max} &= \tfrac{1}{4}\alpha\,(R_c/L_p)^{-2}\,c^4/G \quad
             \text{(unchanged from prior canon)}.
\end{align}
No new free parameters are introduced; the
Compton frequency $\omC$ replaces $\rc$ one-for-one.

% ═════════════════════════════════════════════════════════════════════════════
\section{Numerical Verification}
% ═════════════════════════════════════════════════════════════════════════════

\begin{table}[h]
\centering
\caption{Numerical verification of all Compton-derived constants
  against SST canonical values (CODATA 2022~\cite{MohrTaylor2016}).}
\label{tab:numerics}
\begin{tabular}{lSSS}
  \toprule
  Quantity & {Derived value} & {Canon value} & {Relative error} \\
  \midrule
  $\rc$ / \si{m}
    & 1.40897e-15 & 1.40897017e-15 & {$<10^{-6}$} \\
  $\GammaZ$ / \si{m^2.s^{-1}}
    & 9.6836e-9   & 9.68361918e-9  & {$<10^{-6}$} \\
  $\alpha^{-1}$
    & 137.036     & 137.035999177  & {$6\times10^{-8}$} \\
  $\rhocore$ / \si{kg.m^{-3}}
    & 7.787e18   & 3.89343583e18  & {factor 2 open} \\
  $\mathcal{C}=\rhocore/\rhof$
    & 1.11e25    & 5.562e24       & {factor 2 open} \\
  \bottomrule
\end{tabular}
\end{table}

The sub-part-per-million agreement in Table~\ref{tab:numerics} is
expected: the Compton anchoring
postulate~\eqref{eq:compton_anchor} is by construction consistent
with the definition of the canonical swirl velocity
$\vswirl=\alpha c/2$.  The independent content of the postulate is
that this velocity ratio equals $\omC\rc/c$, which is not a priori
guaranteed unless the vortex boundary rotates at the Compton
frequency.

% ═════════════════════════════════════════════════════════════════════════════
\section{Summary of Resolved and Open Problems}
% ═════════════════════════════════════════════════════════════════════════════

\begin{table}[h]
\centering
\caption{Status of the three Route-2 open problems after the present work.}
\label{tab:status}
\begin{tabular}{clp{7.5cm}l}
  \toprule
  \# & Problem & Resolution & Status \\
  \midrule
  (i)  & $\rc$ is an external constant
       & Derived via $\rc=\vswirl/\omC$,
         equation~\eqref{eq:rc_derived}.
       & \textbf{Resolved} \\[4pt]
  (ii) & $\GammaZ$ defined self-referentially
       & Derived via $\GammaZ=2\pi\vswirl^2/\omC$,
         equation~\eqref{eq:Gamma0_derived}.
       & \textbf{Resolved} \\[4pt]
  (iii)& Energy-scale gap
       & Functional form $\rhocore\propto m_ec^2/(\vswirl^2\rc^3)$
         derived via $F_{\max}$ constraint,
         equation~\eqref{eq:rhocore_final}.
         Geometrical prefactor (factor of two) requires further work.
       & \textit{Substantially advanced; prefactor open} \\
  \bottomrule
\end{tabular}
\end{table}

Problems~(i) and~(ii) are cleanly resolved by the Compton anchoring
postulate.  Problem~(iii) is substantially advanced: the functional
form $\rhocore\propto m_ec^2/(\vswirl^2\rc^3)$ is derived without
circularity, and the energy-scale bridge is established.  A residual
factor of two in the geometrical prefactor of $\rhocore$ requires a
more careful treatment of the vortex-boundary pressure balance and
is left as an explicit research item.

The key epistemic caveat is that $\omC=m_ec^2/\hbar$
introduces $m_e$, $c$, and $\hbar$ as inputs.  Within the current
SST programme these are treated as measured constants~\cite{MohrTaylor2016},
consistent with the Zero-Parameter Principle of the
canon~\cite{IskandaraniCanon078}: no new \emph{free} parameters are
added, only the replacement of two calibrated inputs
($\rc$ and $\rhocore$) by a single calibrated input ($\omC$) that is
more fundamental in the quantum-mechanical hierarchy.  Whether $m_e$
itself is derivable from $(\rhof,\vswirl,\omC)$ is a question for a
future stage of the programme.

% ─────────────────────────────────────────────────────────────────────────────
\begin{thebibliography}{99}

\bibitem{IskandaraniCanon078}
O.~Iskandarani,
\textit{Swirl-String Theory Canon v0.7.8},
Zenodo, 2025.
\href{https://doi.org/10.5281/zenodo.17619170}{DOI:~10.5281/zenodo.17619170}.

\bibitem{IskandaraniSST59}
O.~Iskandarani,
\textit{Mass Generation via Topological Shielding Gates in
Swirl-String Theory (SST-59)},
Zenodo, 2025.
\href{https://doi.org/10.5281/zenodo.17619573}{DOI:~10.5281/zenodo.17619573}.

\bibitem{IskandaraniVAMCore2025}
O.~Iskandarani,
\textit{Core Foundations of the Vortex \AE{}ther Model},
Zenodo, 2025.
\href{https://doi.org/10.5281/zenodo.16424683}{DOI:~10.5281/zenodo.16424683}.

\bibitem{Saffman1992}
P.~G.~Saffman,
\textit{Vortex Dynamics},
Cambridge University Press, 1992.
\href{https://doi.org/10.1017/CBO9780511624063}{DOI:~10.1017/CBO9780511624063}.

\bibitem{Gilbert2016}
B.~Gilbert,
\textit{Database of Ideal Knots, 3--10 Crossings},
2016.
Dataset of Fourier coefficients and ropelength data for ideal knots.

\bibitem{MohrTaylor2016}
P.~J.~Mohr, D.~B.~Newell, B.~N.~Taylor, and E.~Tiesinga,
\textit{2022 CODATA Recommended Values of the Fundamental
Physical Constants},
Rev.\ Mod.\ Phys.\ (2024).
\href{https://physics.nist.gov/constants}{physics.nist.gov/constants}.

\bibitem{Hestenes1990zitterbewegung}
D.~Hestenes,
``The zitterbewegung interpretation of quantum mechanics,''
Found.\ Phys.\ \textbf{20}, 1213--1232 (1990).
\href{https://doi.org/10.1007/BF01889466}{DOI:~10.1007/BF01889466}.

\bibitem{Barut1981}
A.~O.~Barut and A.~J.~Bracken,
``Zitterbewegung and the internal geometry of the electron,''
Phys.\ Rev.\ D \textbf{23}, 2454--2461 (1981).
\href{https://doi.org/10.1103/PhysRevD.23.2454}{DOI:~10.1103/PhysRevD.23.2454}.

\bibitem{Donnelly1991}
R.~J.~Donnelly,
\textit{Quantized Vortices in Helium~II},
Cambridge University Press, 1991.

\bibitem{Feynman1955}
R.~P.~Feynman,
``Application of quantum mechanics to liquid helium,''
Prog.\ Low Temp.\ Phys.\ \textbf{1}, 17--53 (1955).
\href{https://doi.org/10.1016/S0079-6417(08)60077-3}{DOI:~10.1016/S0079-6417(08)60077-3}.

\bibitem{ChatGPTAudit2025}
ChatGPT (OpenAI),
\textit{Technical audit of the SST Route-2 robustness sweep},
private communication, 2025.

\bibitem{LandauMechanics}
L.~D.~Landau and E.~M.~Lifshitz,
\textit{Mechanics}, 3rd ed.,
Butterworth-Heinemann, 1976, \S10 (virial theorem).

\end{thebibliography}

\end{document}
